\documentclass[preprint,3p,number]{elsarticle}

\usepackage{amsmath,amssymb}

\title{Derivation of Dimentionless Governing Equations for Two-Dimensional Natural Convection}
\author{}
\date{}

\begin{document}

\maketitle

\section{Derivation of Dimensionless Weak-Form Governing Equations}

For a two-dimensional, steady-state, incompressible pure-water
natural convection model, the dimensionless strong-form governing
equations are given by

\subsection{Dimensionless Strong-Form Governing Equations}

\textbf{Continuity equation:}
\begin{equation}
\frac{\partial U}{\partial X}
+\frac{\partial V}{\partial Y}=0
\end{equation}

\textbf{Horizontal momentum equation:}
\begin{equation}
U\frac{\partial U}{\partial X}
+V\frac{\partial U}{\partial Y}
=-\frac{\partial P}{\partial X}
+\mathrm{Pr}\left(
\frac{\partial^2 U}{\partial X^2}
+\frac{\partial^2 U}{\partial Y^2}
\right)
\end{equation}

\textbf{Vertical momentum equation:}
\begin{equation}
U\frac{\partial V}{\partial X}
+V\frac{\partial V}{\partial Y}
=-\frac{\partial P}{\partial Y}
+\mathrm{Pr}\left(
\frac{\partial^2 V}{\partial X^2}
+\frac{\partial^2 V}{\partial Y^2}
\right)
+\mathrm{Ra}\,\mathrm{Pr}\,\theta
\end{equation}

\textbf{Energy equation:}
\begin{equation}
U\frac{\partial \theta}{\partial X}
+V\frac{\partial \theta}{\partial Y}
=\frac{\partial^2\theta}{\partial X^2}
+\frac{\partial^2\theta}{\partial Y^2}
\end{equation}

Let $\Omega$ denote the computational domain and
$\partial\Omega$ its boundary. The test functions for the
continuity, horizontal momentum, vertical momentum and energy
equations are denoted by $q$, $w_x$, $w_y$ and $\phi$,
respectively.

\subsection{Weak Form of the Continuity Equation}

Multiplying the continuity equation by the test function $q$
and integrating over the domain $\Omega$ gives

\begin{equation}
\int_{\Omega}
q\left(
\frac{\partial U}{\partial X}
+\frac{\partial V}{\partial Y}
\right)\,d\Omega=0
\end{equation}

Therefore, the dimensionless weak form of the continuity equation is

\begin{equation}
\boxed{
\int_{\Omega}q\,\nabla\cdot\mathbf{U}\,d\Omega=0
}
\end{equation}

where $\mathbf{U}=(U,V)$ is the dimensionless velocity vector.

\subsection{Weak Form of the Momentum Equation}

The horizontal and vertical momentum equations can be combined
into the following dimensionless vector equation:

\begin{equation}
(\mathbf{U}\cdot\nabla)\mathbf{U}
=-\nabla P
+\mathrm{Pr}\nabla^2\mathbf{U}
+\mathrm{Ra}\,\mathrm{Pr}\,\theta\mathbf{e}_Y
\end{equation}

where $\mathbf{U}=(U,V)$ is the dimensionless velocity vector
and $\mathbf{e}_Y=(0,1)^T$ is the unit vector in the positive
vertical direction.

Let $\mathbf{W}=(w_x,w_y)$ be the vector test function.
Multiplying the momentum equation by $\mathbf{W}$ using the
dot product and integrating over the domain $\Omega$ gives

\begin{align}
&\int_{\Omega}\mathbf{W}\cdot
[(\mathbf{U}\cdot\nabla)\mathbf{U}]\,d\Omega
\nonumber\\
&=-\int_{\Omega}\mathbf{W}\cdot\nabla P\,d\Omega
+\mathrm{Pr}\int_{\Omega}
\mathbf{W}\cdot\nabla^2\mathbf{U}\,d\Omega
\nonumber\\
&\quad+\mathrm{Ra}\,\mathrm{Pr}
\int_{\Omega}\theta W_Y\,d\Omega
\end{align}

Applying integration by parts to the pressure term gives

\begin{equation}
-\int_{\Omega}\mathbf{W}\cdot\nabla P\,d\Omega
=
-\int_{\partial\Omega}
P(\mathbf{W}\cdot\mathbf{n})\,d\Gamma
+\int_{\Omega}P(\nabla\cdot\mathbf{W})\,d\Omega
\end{equation}

Similarly, integration by parts of the viscous term gives

\begin{align}
\mathrm{Pr}\int_{\Omega}
\mathbf{W}\cdot\nabla^2\mathbf{U}\,d\Omega
&=-\mathrm{Pr}\int_{\Omega}
\nabla\mathbf{W}:\nabla\mathbf{U}\,d\Omega
\nonumber\\
&\quad+\mathrm{Pr}\int_{\partial\Omega}
\mathbf{W}\cdot\frac{\partial\mathbf{U}}{\partial n}
\,d\Gamma
\end{align}

Substituting these expressions and rearranging yields

\begin{align}
&\int_{\Omega}\mathbf{W}\cdot
[(\mathbf{U}\cdot\nabla)\mathbf{U}]\,d\Omega
\nonumber\\
&+\mathrm{Pr}\int_{\Omega}
\nabla\mathbf{W}:\nabla\mathbf{U}\,d\Omega
\nonumber\\
&-\int_{\Omega}P(\nabla\cdot\mathbf{W})\,d\Omega
\nonumber\\
&-\mathrm{Ra}\,\mathrm{Pr}
\int_{\Omega}\theta W_Y\,d\Omega
=B_m
\end{align}

where the boundary contribution is

\begin{align}
B_m={}&-\int_{\partial\Omega}
P(\mathbf{W}\cdot\mathbf{n})\,d\Gamma
\nonumber\\
&+\mathrm{Pr}\int_{\partial\Omega}
\mathbf{W}\cdot\frac{\partial\mathbf{U}}{\partial n}
\,d\Gamma
\end{align}

For homogeneous no-slip velocity boundary conditions,
$\mathbf{W}=\mathbf{0}$ on $\partial\Omega$, and therefore
$B_m=0$.

Hence, the final dimensionless weak form of the momentum equation
is

\begin{equation}
\boxed{
\begin{aligned}
&\int_{\Omega}\mathbf{W}\cdot
[(\mathbf{U}\cdot\nabla)\mathbf{U}]\,d\Omega\\
&+\mathrm{Pr}\int_{\Omega}
\nabla\mathbf{W}:\nabla\mathbf{U}\,d\Omega\\
&-\int_{\Omega}P(\nabla\cdot\mathbf{W})\,d\Omega\\
&-\mathrm{Ra}\,\mathrm{Pr}
\int_{\Omega}\theta W_Y\,d\Omega=0
\end{aligned}
}
\end{equation}

\subsection{Weak Form of the Energy Equation}

Multiplying the energy equation by the test function $\phi$
and integrating over $\Omega$ gives

\begin{align}
&\int_{\Omega}\phi
\left(
U\frac{\partial\theta}{\partial X}
+V\frac{\partial\theta}{\partial Y}
\right)d\Omega \nonumber\\
&=\int_{\Omega}\phi
\left(
\frac{\partial^2\theta}{\partial X^2}
+\frac{\partial^2\theta}{\partial Y^2}
\right)d\Omega
\end{align}

Applying integration by parts to the diffusion term gives

\begin{align}
&\int_{\Omega}\phi
\left(
\frac{\partial^2\theta}{\partial X^2}
+\frac{\partial^2\theta}{\partial Y^2}
\right)d\Omega \nonumber\\
&=-\int_{\Omega}
\left(
\frac{\partial\phi}{\partial X}
\frac{\partial\theta}{\partial X}
+\frac{\partial\phi}{\partial Y}
\frac{\partial\theta}{\partial Y}
\right)d\Omega \nonumber\\
&\quad+\int_{\partial\Omega}
\phi\frac{\partial\theta}{\partial n}\,d\Gamma
\end{align}

Rearranging gives

\begin{align}
&\int_{\Omega}\phi
\left(
U\frac{\partial\theta}{\partial X}
+V\frac{\partial\theta}{\partial Y}
\right)d\Omega \nonumber\\
&+\int_{\Omega}\nabla\phi\cdot\nabla\theta\,d\Omega
=B_\theta
\end{align}

where the boundary contribution is

\begin{equation}
B_\theta=
\int_{\partial\Omega}
\phi\frac{\partial\theta}{\partial n}\,d\Gamma
\end{equation}

For the prescribed-temperature vertical walls, the test function
$\phi$ vanishes on the boundary. For the insulated top and bottom
walls, $\partial\theta/\partial n=0$. Therefore, $B_\theta=0$,
and the dimensionless weak form of the energy equation becomes

\begin{equation}
\boxed{
\begin{aligned}
&\int_{\Omega}\phi
\left(
U\frac{\partial\theta}{\partial X}
+V\frac{\partial\theta}{\partial Y}
\right)d\Omega\\
&+\int_{\Omega}\nabla\phi\cdot\nabla\theta\,d\Omega=0
\end{aligned}
}
\end{equation}

\subsection{Summary of the Dimensionless Weak Forms}

The final dimensionless weak-form governing equations for the
two-dimensional pure-water natural convection model are:

\textbf{Continuity:}
\begin{equation}
\int_{\Omega}q\,\nabla\cdot\mathbf{U}\,d\Omega=0
\end{equation}

\textbf{Momentum:}
\begin{equation}
\begin{aligned}
&\int_{\Omega}\mathbf{W}\cdot
[(\mathbf{U}\cdot\nabla)\mathbf{U}]\,d\Omega\\
&+\mathrm{Pr}\int_{\Omega}
\nabla\mathbf{W}:\nabla\mathbf{U}\,d\Omega\\
&-\int_{\Omega}P(\nabla\cdot\mathbf{W})\,d\Omega\\
&-\mathrm{Ra}\,\mathrm{Pr}
\int_{\Omega}\theta W_Y\,d\Omega=0
\end{aligned}
\end{equation}

\textbf{Energy:}
\begin{equation}
\begin{aligned}
&\int_{\Omega}\phi
\left(
U\frac{\partial\theta}{\partial X}
+V\frac{\partial\theta}{\partial Y}
\right)d\Omega\\
&+\int_{\Omega}\nabla\phi\cdot\nabla\theta\,d\Omega=0
\end{aligned}
\end{equation}

\end{document}
